\subsection{Exact small-momentum dispersion}
\label{sec:lovelock-momentum}

Let \(k\) be the magnitude of the spatial wave vector and
\(w=\omega/2\). With the same curvature coupling and mass
parameter \(\nu\) as in section~\ref{sec:detuning}, the reduced
equation at finite \(k\) is
\begin{equation}
 (rfA')'+\left(\frac{\omega^2 r}{f}
      -(\nu^2-1)r-\frac{k^2}{r}\right)A=0 .
 \label{eq:lovelock-k-radial}
\end{equation}
In \eqref{eq:detune-v} this adds the potential
\(-k^2/[4v(v^2-1)]\), with a fourth regular singular point at
\(v=0\). The zero-momentum solution suffices to find the exact
first derivative of every quasinormal pole.

Use \(N=n+h\), \(h=(1+\nu)/2\), for the overtone \(n\geq0\).
On the positive-real-part branch the unperturbed pole is
\(\omega_n^0=2(1-\ii)N\). It is simple because only the gamma
factor with argument \(h-s_+\) gives a zero of the source there.
Write \(x=(v-1)/(v+1)\) and \(\alpha=-\ii\omega_n^0/2\).
The polynomial at the pole can be converted by Pfaff's transformation
to a Jacobi polynomial normalized to unity at \(x=0\):
\begin{align}
 A_n(x)&=x^{\alpha/2}(1-x)^h \mathcal P_n(x),\notag\\
 \mathcal P_n(x)&={}_2F_1(-n,n+\alpha+\nu+1;\alpha+1;x).
 \label{eq:lovelock-jacobi-mode}
\end{align}
The overall normalization is arbitrary. The endpoint \(x=0\)
is the horizon and \(x=1\) is the AdS boundary.

Set \(b_n=\partial_{k^2}\omega_n|_{k=0}\). Differentiate
\eqref{eq:lovelock-k-radial} along the pole, multiply by
\(A_n\), and subtract the equation for \(A_n\) times the
differentiated mode. Integrating the resulting Wronskian gives
\begin{equation}
 2\omega_n^0 b_n \operatorname{FP}\!\int_1^\infty
      \frac{r}{f}A_n^2\,\dd r
   =\operatorname{FP}\!\int_1^\infty \frac{A_n^2}{r}\,\dd r .
 \label{eq:lovelock-bilinear}
\end{equation}
Here \(\operatorname{FP}\) denotes the constant part after subtracting
the horizon power terms. The boundary Wronskian vanishes because
both modes have zero source. At the horizon the variation of the
ingoing exponent produces a logarithm in the differentiated mode.
In the Wronskian the logarithm cancels, leaving powers
\(x^{\alpha+j}\) with integer \(j\geq0\). These terms have vanishing
finite part. Equivalently, the moment integrals in \(x\) are analytic
continuations from \(\operatorname{Re}\alpha>0\). Equation
\eqref{eq:lovelock-bilinear} is bilinear in the mode.

With weight \(x^\alpha(1-x)^\nu\), the two integrals use
\((1+x)/(4x)\) and \(1/(1+x)\), respectively. Jacobi orthogonality
replaces \(\mathcal P_n^2/(1+x)\) in the second integral by
\(\mathcal P_n(-1)\mathcal P_n(x)/(1+x)\), since the difference
contains a polynomial of degree \(n-1\) times \(\mathcal P_n\).
The same argument replaces \(\mathcal P_n^2/x\) by
\(\mathcal P_n/x\). Integrating the Jacobi Rodrigues formula \cite{NIST:DLMF}
by parts then gives
\begin{equation}
 b_n=
 \frac{2\mathcal P_n(-1)}
 {\omega_n^0\left(1+\dfrac{\alpha}{2n+\alpha+\nu+1}\right)}
 \frac{(\alpha)_{n+1}}{(\alpha+\nu+n+1)_{n+1}}\,
 {}_2F_1(n+1,\alpha+n+1;\alpha+\nu+2n+2;-1).
 \label{eq:lovelock-exact-dispersion}
\end{equation}
To see the normalization, put
\(\mathcal I_{-1}=n!B(\alpha,n+\nu+1)/(\alpha+1)_n\) for the
continued \(1/x\) moment. The Jacobi squared norm is
\(\mathcal I_{-1}\alpha/(2n+\alpha+\nu+1)\).
The continued \(1/(1+x)\) moment is the product of
\(\mathcal P_n(-1)n!B(\alpha+n+1,n+\nu+1)/(\alpha+1)_n\)
and the hypergeometric function in \eqref{eq:lovelock-exact-dispersion}.
The beta function ratio produces the two Pochhammer symbols.
The negative-real-part branch obeys
\(\omega_n^-(k)=-[\omega_n^+(k)]^*\) for real \(k\).

\begin{table}[t]
 \centering\small
 \begin{tabular}{cccc}
 \hline
 \(\nu\) & \(n\) & \(b_n\) from \eqref{eq:lovelock-exact-dispersion}
     & Absolute error from pole matching \\
 \hline
 0.75 & 0 & $0.45885866+0.13218415\ii$ & 2.73e-10 \\
0.75 & 3 & $0.17369969+0.04694990\ii$ & 2.20e-12 \\
1 & 5 & $0.12944937+0.03540774\ii$ & 5.00e-13 \\
1.25 & 0 & $0.36222059+0.13355222\ii$ & 5.80e-11 \\
1.25 & 8 & $0.09782584+0.02683118\ii$ & 6.88e-14 \\
\hline

 \end{tabular}
 \caption{First momentum derivative of the Lovelock scalar poles.
 The numerical column reports the difference between the analytic
 result and independent finite-momentum pole matching.}
 \label{tab:lovelock-momentum}
\end{table}

For table~\ref{tab:lovelock-momentum} we match the ingoing and
normalizable Frobenius series in \(u=r^{-2}\) for the full
equation \eqref{eq:lovelock-k-radial}. We use two centered
differences in \(k^2\), with steps \(0.02\) and \(0.01\), and
extrapolate quadratically in the step. The calculation also
changes the matching point from \(u=0.54\) to \(u=0.61\) and
increases the number of series terms.

\subsection{Airy limit of the momentum shift}
\label{sec:lovelock-airy}

The large-overtone momentum scale follows from the singularity
at \(z_s=\ii\pi/2\) on the upper optical-time branch. The equation
at finite \(k\) in the coordinate of \eqref{eq:lovelock-z} is
\begin{equation}
 A_{zz}+\left[
  \frac{\omega^2-k^2}{4}\coth z+\frac{k^2}{4}\tanh z
        -\frac{\nu^2-1}{4\sinh^2 z}\right]A=0 .
 \label{eq:lovelock-k-z}
\end{equation}
For \(y=w^{2/3}(z-z_s)\), the leading inner equation is
\(A_{yy}+[y+\lambda/y]A=0\), with
\(\lambda=k^2/(4w^{2/3})\). This follows from the simple zero of
\(\coth z\) and the simple pole of \(\tanh z\) at \(z_s\).
The terms depending on the bare mass vanish at leading order
when \(\nu\) is fixed.

For the matching phase, let \(\mathcal Z_{\rm L}=\pi(1+\ii)/4\)
be the optical distance on the branch reached below the
physical horizon. The boundary ray \(r=e^{-\ii\pi/4}R\), \(R>0\),
has \(x_{\rm opt}-\mathcal Z_{\rm L}\) of phase \(-3\pi/4\)
near the singularity. Here \(x_{\rm opt}\) is the tortoise
coordinate with zero at the boundary. At the unperturbed pole
\(\arg w=-\pi/4\), the Airy coordinate \(X=e^{\ii\pi/3}y\)
has phase \(-\pi/3\) on that boundary ray, since
\(z-z_s\sim r^2\). The adjacent ray towards the physical
horizon has \(\arg r=\pi/12\) near the core, giving
\(\arg X=\pi/3\). In \(X\) the inner equation becomes
\begin{equation}
 F''-XF+\frac{g}{X}F=0,\qquad
 g=e^{-\ii\pi/3}\frac{k^2}{4w^{2/3}}.
 \label{eq:lovelock-perturbed-airy}
\end{equation}
Across the two matching rays the continuation of the
ingoing exponential changes by the Stokes multiplier of
the outgoing solution, recessive on \(X>0\).

Normalize the recessive solution \(F(X;g)\) to \(\operatorname{Ai}(X)\)
as \(X\to+\infty\), and let \(a(g)=F(0;g)\).
Although the derivative can have a logarithm at the origin,
\(a(g)\) is finite. With \(\sigma=e^{2\pi\ii/3}\), the rotated
solutions are \(F_j(X)=F(\sigma^{-j}X;\sigma^j g)\).
The WKB-normalized adjacent Wronskians have their values at \(g=0\),
since the \(g/X\) term makes a correction that tends to zero
in the large-\(X\) asymptotics. Thus
\(F_{-1}+\sigma F_1=S_{\rm A}(g)F_0\).
Evaluating this relation at \(X=0\) gives
\begin{equation}
 S_{\rm A}(0)=-\sigma^{-1},\qquad
 \left.\partial_g\log S_{\rm A}\right|_{0}=-3\,\frac{a'(0)}{a(0)}.
\end{equation}

Variation of parameters with \(\operatorname{Ai}\) and
\(\operatorname{Bi}\) yields the convergent integral
\begin{equation}
 \frac{a'(0)}{a(0)}
 =\pi\int_0^\infty\frac{\operatorname{Ai}(t)
 [\sqrt3\,\operatorname{Ai}(t)-\operatorname{Bi}(t)]}{t}\,\dd t
 =-\frac{\Gamma(1/3)^2}{3^{4/3}}.
 \label{eq:lovelock-airy-moment}
\end{equation}
For the last equality, put \(s=2t^{3/2}/3\) and use the Airy
representations in terms of \(I_{1/3}\) and \(K_{1/3}\) \cite{NIST:DLMF}.
The resulting Bessel moment is
\begin{equation}
 \int_0^\infty s^{-1/3} I_{1/3}(s)K_{1/3}(s)\,\dd s
 =2^{-4/3}\Gamma(1/3)^2.
\end{equation}
It follows by Mellin transformation of the product representation
\begin{equation}
 I_{1/3}(s)K_{1/3}(s)=\int_0^\infty
      e^{-2v/3}J_0(2s\sinh v)\,\dd v
\end{equation}
and continuation from the strip of absolute convergence.

The leading pole condition uses the boundary phase
\(e^{-\ii\pi\nu}\) and the round-trip symbol
\(e^{2\ii\omega\mathcal Z_{\rm L}}\). Its reflection coefficient
is fixed to \(-1\) at \(g=0\) by the exact pole locations.
The Airy matching multiplies it by \(S_{\rm A}(g)/S_{\rm A}(0)\),
since the ratio of the WKB amplitudes changes by the same
Stokes coefficient. Exterior momentum corrections from outside
a matching radius \(|z-z_s|=\varepsilon\) are of order
\(k^2 |w|^{-1}\varepsilon^{-1/2}\), which is subleading when
\(\varepsilon\gg |w|^{-2/3}\) and \(\varepsilon\to0\).
The resulting leading WKB prediction for the slope is
\begin{align}
 N^{2/3} b_n &\longrightarrow
 \mathcal B_{\rm A}
 =\frac{\ii e^{-\ii\pi/3}\Gamma(1/3)^2}
        {8\,\mathcal Z_{\rm L}\,3^{1/3}(1-\ii)^{2/3}}\notag\\
 &=0.42932943+0.11503847\ii .
 \label{eq:lovelock-airy-limit}
\end{align}
The power of \(1-\ii\) is continued from real positive \(w\).
Equation \eqref{eq:lovelock-exact-dispersion} gives an independent
test of the complex coefficient in this prediction.

\begin{table}[t]
 \centering\small
 \begin{tabular}{cccc}
 \hline
 \(\nu\) & \(n=16\) & \(n=64\) & \(n=256\) \\
 \hline
 0.75 & 1.579e-04 & 2.622e-05 & 4.185e-06 \\
1.25 & 1.309e-03 & 2.208e-04 & 3.539e-05 \\
2 & 3.993e-03 & 6.897e-04 & 1.112e-04 \\
\hline

 \end{tabular}
 \caption{Absolute error \(\lvert N^{2/3}b_n-\mathcal B_{\rm A}\rvert\)
 using the exact momentum slopes in \eqref{eq:lovelock-exact-dispersion}.}
 \label{tab:lovelock-airy}
\end{table}

Table~\ref{tab:lovelock-airy} tests the coefficient at several
values of the bare mass. The scale \(k\sim|\omega|^{1/3}\)
controls the Airy reflection equation at fixed \(\nu\).
Because the Lovelock lapse has the same radial form as in AdS\(_5\),
the pure optical geodesics here have
\(J_{n,0}=n\pi-\arctan(2/E^2)\) from \eqref{eq:jacobi-ads5},
with transverse multiplicity \(d-1\).
The worldline Gaussian factor assumes a large mass and a transverse
Gaussian width small compared to the inner turning radius.
The following calculation connects the momentum derivative in the
two regimes.

\subsection{Joint mass and overtone limit}
\label{sec:lovelock-crossover}

The Airy equation also resolves the regime in which the mass
term in \eqref{eq:lovelock-k-z} remains finite in the inner scaling.
Let \(N=n+h\) as above, take \(\nu\geq1\), and hold
\begin{equation}
 \chi=\frac{\sqrt{\nu^2-1}}{N^{2/3}}
 \label{eq:lovelock-scaled-mass}
\end{equation}
fixed as \(N\to\infty\). We compute the slope at \(k^2=0\).
Since \(\sinh^2 z_s=-1\), the mass contributes a constant
\((\nu^2-1)/(4w^{4/3})\) to the leading inner equation in \(y\).
In the coordinate \(X=e^{\ii\pi/3}y\) of
\eqref{eq:lovelock-perturbed-airy}, the equation becomes
\begin{equation}
 F''-\left(X-\zeta-\frac gX\right) F=0,\qquad
 \zeta=e^{-\ii\pi/3}\frac{\chi^2}{4\,2^{2/3}}.
 \label{eq:lovelock-displaced-airy}
\end{equation}
Here \(w=(1-\ii)N\) at \(k=0\), and \(g\) has the same
definition as in \eqref{eq:lovelock-perturbed-airy}.
The mass displaces the Airy turning point. The momentum
perturbation remains anchored at the singularity \(X=0\).

\subsubsection{The connection coefficient at fixed scaled mass}

Normalize \(F(X;g,\zeta)\) to \(\operatorname{Ai}(X-\zeta)\)
at positive large \(X\), and define
\(a(g,\zeta)=F(0;g,\zeta)\) and
\(a_1(\zeta)=\partial_g a(0,\zeta)\).
Variation of parameters gives
\begin{multline}
 a_1(\zeta)=\pi\int_0^\infty\frac{\operatorname{Ai}(t-\zeta)}t
 \big[\operatorname{Bi}(-\zeta)\operatorname{Ai}(t-\zeta)\\
       -\operatorname{Ai}(-\zeta)\operatorname{Bi}(t-\zeta)\big]\dd t .
 \label{eq:lovelock-shifted-variation}
\end{multline}
The bracket vanishes linearly at the origin, and the product
\(\operatorname{Ai}\operatorname{Bi}\) decays as \(t^{-1/2}\)
at large positive \(t\). With \(\sigma=e^{2\pi\ii/3}\), the
rotated solutions now use
\((X,g,\zeta)\mapsto(\sigma^{-j}X,\sigma^j g,\sigma^{-j}\zeta)\).
Each adjacent Wronskian keeps its \(g=0\) value by the
same large-\(X\) normalization as in section~\ref{sec:lovelock-airy}.
At \(g=0\) the connection relation is the Airy identity
with argument \(X-\zeta\), so \(S_{\rm A}(0,\zeta)=-\sigma^{-1}\).
Taking its \(g\) derivative at the origin gives
\begin{equation}
 D_{\rm A}(\zeta)\equiv\partial_g\log S_{\rm A}(0,\zeta)
 =-\frac{a_1(\zeta)+a_1(\sigma\zeta)+a_1(\sigma^{-1}\zeta)}
         {\operatorname{Ai}(-\zeta)} .
 \label{eq:lovelock-displaced-connection}
\end{equation}
At \(\zeta=0\) this equals \(\Gamma(1/3)^2/3^{1/3}\), recovering
\eqref{eq:lovelock-airy-limit}.

At fixed \(\chi\), the shifted Airy region has size
\(|w|^{-2/3}\) in \(z\). The exterior momentum derivative
is subleading by the same overlap estimate as in the
fixed-mass calculation, since the mass shift in \(z\) is
of this order as well. The matched pole condition
therefore predicts
\begin{equation}
 N^{2/3} b_n\ \longrightarrow\
 \mathcal B(\chi)=
 \frac{\ii e^{-\ii\pi/3}}{8\mathcal Z_{\rm L}(1-\ii)^{2/3}}
 D_{\rm A}(\zeta).
 \label{eq:lovelock-crossover}
\end{equation}
Only the \(k^2\) coefficient of the pole condition is varied
here. Its value at zero momentum follows from the exact spectrum.

\begin{table}[t]
 \centering\small
 \begin{tabular}{cccc}
 \hline
 \(\chi\) & \(\mathcal B(\chi)\) & Error at \(n=64\) & Error at \(n=256\) \\
 \hline
 0 & $0.4293294+0.1150385\ii$ & 1.12e-04 & 1.78e-05 \\
1 & $0.3922517+0.1456450\ii$ & 8.31e-05 & 1.40e-05 \\
2 & $0.2873197+0.1840777\ii$ & 5.25e-05 & 9.50e-06 \\
3 & $0.1811542+0.1651296\ii$ & 5.23e-05 & 1.01e-05 \\
5 & $0.1002772+0.1002742\ii$ & 9.90e-06 & 2.24e-06 \\
\hline

 \end{tabular}
 \caption{Joint mass and overtone limit of the momentum slopes.
 The errors compare \eqref{eq:lovelock-crossover} with
 \(N^{2/3}b_n\) computed from \eqref{eq:lovelock-exact-dispersion}
 at the indicated fixed \(\chi\).}
 \label{tab:lovelock-crossover}
\end{table}

Table~\ref{tab:lovelock-crossover} compares the local connection
calculation with the exact moments of the boundary-value problem.
For each \(n\), the mass parameter in the exact formula is chosen
to satisfy \eqref{eq:lovelock-scaled-mass}. In the Airy integrals
we use the Maclaurin series on a finite interval and integrate
the asymptotic expansion of \(\operatorname{Ai}\operatorname{Bi}\)
in the tail \cite{NIST:DLMF}. Changing the matching radius
from 13 to 16 and increasing the truncation order checks that
this numerical error is below the finite-overtone effect.

\subsubsection{The geometric limit of the momentum derivative}

Take \(\chi\) large after the joint limit. In \(X/\zeta\) the
simple-turning-point problem has large WKB parameter
\(\zeta^{3/2}\), and the distant pole at \(X=0\) stays outside
the connection from the outer oscillatory region to the
displaced turning point. The leading \(g\) derivative of
the WKB action from the turning point to infinity is
\begin{equation}
 -\frac12\int_0^\infty\frac{\dd t}{(\zeta+t)\sqrt t}
    =-\frac{\pi}{2\sqrt\zeta}.
 \label{eq:lovelock-displaced-action}
\end{equation}
The endpoint derivative of the action vanishes at the
turning point. The path \(X=\zeta+t\) and the square root
are continued from positive \(\zeta\) to the optical branch.
The simple-turning-point multiplier at the root has its universal
Airy value. Expressing the incoming and outgoing modes in the
normalization at infinity multiplies their ratio by the exponential
of minus twice the action shift in \eqref{eq:lovelock-displaced-action}.
Thus the leading WKB term is \(D_{\rm A}(\zeta)\sim\pi/\sqrt\zeta\),
and \eqref{eq:lovelock-crossover} gives
\begin{equation}
 \mathcal B(\chi)\sim\frac{1+\ii}{2\chi}.
 \label{eq:lovelock-massive-limit}
\end{equation}
For example, at \(\chi=5\) the Airy quadrature gives
\(\chi\mathcal B(\chi)=0.50139+0.50137\ii\).
For \(\chi=10\) and \(n=256\), the exact pole slope gives
\(\chi N^{2/3}b_n\simeq0.500019(1+\ii)\).

The same momentum derivative follows from the optical geodesic.
In this metric the inner-to-outer cycle has
\(J_{\rm opt}=\pi\), as in \eqref{eq:jacobi-ads5}.
If \(M\simeq\nu\) is the large worldline mass and \(k=MP\),
then the optical cycle action satisfies
\(\partial_{P^2}W_{\rm opt}=-J_{\rm opt}/2\).
Consequently the logarithmic \(k^2\) derivative of
\(\exp(-MW_{\rm opt})\) equals \(\pi/(2M)\). In the
pole condition with optical phase \(2\omega\mathcal Z_{\rm L}\)
this yields \(b_n\sim\ii\pi/(4M\mathcal Z_{\rm L})=(1+\ii)/(2M)\),
in agreement with \eqref{eq:lovelock-massive-limit}.
For \(|E|=|\omega|/M\), the condition \(M|E|^{-2}\gg1\) for the
geodesic approximation is \(\chi^3/8\gg1\) in this limit.
Each additional optical reflection adds the same quadratic
momentum term in the exponent. Its Gaussian integral has the
transverse multiplicity \(d-1\) and the repetition dependence
of \eqref{eq:jacobi-ratio}. The comparison fixes this part of the
amplitude in the overlap regime. The full correlator also requires
the residues and contour weights. Section~\ref{sec:lovelock-residues}
determines the residue shifts at first order in \(k^2\).

